%%
%% Copyright 2007-2024 Elsevier Ltd
%%
%% This file is part of the 'Elsarticle Bundle'.
%% ---------------------------------------------
%%
%% It may be distributed under the conditions of the LaTeX Project Public
%% License, either version 1.3 of this license or (at your option) any
%% later version.  The latest version of this license is in
%%    http://www.latex-project.org/lppl.txt
%% and version 1.3 or later is part of all distributions of LaTeX
%% version 1999/12/01 or later.
%%
%% The list of all files belonging to the 'Elsarticle Bundle' is
%% given in the file `manifest.txt'.
%%
%% Template article for Elsevier's document class `elsarticle'
%% with numbered style bibliographic references
%% SP 2008/03/01
%% $Id: elsarticle-template-num.tex 249 2024-04-06 10:51:24Z rishi $
%%
\documentclass[preprint,12pt]{elsarticle}

%% Use the option review to obtain double line spacing
%% \documentclass[authoryear,preprint,review,12pt]{elsarticle}

%% Use the options 1p,twocolumn; 3p; 3p,twocolumn; 5p; or 5p,twocolumn
%% for a journal layout:
%% \documentclass[final,1p,times]{elsarticle}
%% \documentclass[final,1p,times,twocolumn]{elsarticle}
%% \documentclass[final,3p,times]{elsarticle}
%% \documentclass[final,3p,times,twocolumn]{elsarticle}
%% \documentclass[final,5p,times]{elsarticle}
%% \documentclass[final,5p,times,twocolumn]{elsarticle}

%% For including figures, graphicx.sty has been loaded in
%% elsarticle.cls. If you prefer to use the old commands
%% please give \usepackage{epsfig}

%% The amssymb package provides various useful mathematical symbols
\usepackage{amssymb}
%% The amsmath package provides various useful equation environments.
\usepackage{amsmath}
%% The amsthm package provides extended theorem environments
%% \usepackage{amsthm}

%% The lineno packages adds line numbers. Start line numbering with
%% \begin{linenumbers}, end it with \end{linenumbers}. Or switch it on
%% for the whole article with \linenumbers.
%% \usepackage{lineno}

\journal{Aerospace Science and Technology}

\begin{document}

\begin{frontmatter}

%% Title, authors and addresses

%% use the tnoteref command within \title for footnotes;
%% use the tnotetext command for theassociated footnote;
%% use the fnref command within \author or \affiliation for footnotes;
%% use the fntext command for theassociated footnote;
%% use the corref command within \author for corresponding author footnotes;
%% use the cortext command for theassociated footnote;
%% use the ead command for the email address,
%% and the form \ead[url] for the home page:
%% \title{Title\tnoteref{label1}}
%% \tnotetext[label1]{}
%% \author{Name\corref{cor1}\fnref{label2}}
%% \ead{email address}
%% \ead[url]{home page}
%% \fntext[label2]{}
%% \cortext[cor1]{}
%% \affiliation{organization={},
%%             addressline={},
%%             city={},
%%             postcode={},
%%             state={},
%%             country={}}
%% \fntext[label3]{}

\title{Dynamically feasible risk-aware trajectory planning for fixed-wing UAVs via multi-objective two-stage search}

%% use optional labels to link authors explicitly to addresses:
%% \author[label1,label2]{}
%% \affiliation[label1]{organization={},
%%             addressline={},
%%             city={},
%%             postcode={},
%%             state={},
%%             country={}}
%%
%% \affiliation[label2]{organization={},
%%             addressline={},
%%             city={},
%%             postcode={},
%%             state={},
%%             country={}}

\author{Zhiwen Lin}
\affiliation{organization={},%Department and Organization
            addressline={},
            city={Changsha},
            postcode={410073},
            state={Hunan},
            country={China}
            }

%% Abstract
\begin{abstract}

Abstract text.
\end{abstract}


\begin{graphicalabstract}
\end{graphicalabstract}

\begin{highlights}
\item Research highlight 1
\item Research highlight 2
\end{highlights}

%% Keywords
\begin{keyword}
Fixed-wing UAV \sep
trajectory planning \sep
risk-aware planning \sep
multi-objective optimization \sep
A* search \sep
B-spline

\end{keyword}

\end{frontmatter}
\section{Introduction}
\label{sec:introduction}


% Replace the existing Section 2, including its heading, with this file.
% Requires amsmath and amssymb. All pr: labels are local to this replacement.
\section{Problem formulation}
\label{pr:sec:problem}

\subsection{Three-dimensional mission and threat environment}
\label{pr:sec:environment}
The mission is defined in a three-dimensional Cartesian workspace
$\mathcal{W}\subset\mathbb{R}^{3}$. The UAV position is
$\mathbf{p}=[x,y,z]^{\mathsf T}$, where $x$ and $y$ are horizontal
coordinates and $z$ is altitude. The prescribed mission endpoints are
$\mathbf{p}_{s}$ and $\mathbf{p}_{g}$. A discrete search path is written as
\begin{equation}
 \Gamma_d=(\mathbf{p}_0,\ldots,\mathbf{p}_N),\qquad
 \mathbf{p}_0=\mathbf{p}_{s},\quad \mathbf{p}_N=\mathbf{p}_{g},
 \label{pr:eq:discrete}
\end{equation}
whereas the final time-parameterized trajectory is
$\mathbf{r}(t)\in\mathcal{W}$, $t\in[0,T]$. These representations are
distinguished because collision-free graph connectivity does not establish
continuous-trajectory feasibility. The search, geometric refinement and
speed-assignment procedures are described in Section~3.

Hard exclusion regions and traversable soft-threat regions are modeled
separately. Each hard zone is a vertical cylinder,
\begin{equation}
 \Omega_{F,j}=\left\{\mathbf{p}\in\mathcal{W}:
 (x-x_{F,j})^2+(y-y_{F,j})^2\le R_{F,j}^{2},\quad
 0\le z\le H_{F,j}\right\},
 \label{pr:eq:hardzone}
\end{equation}
where $R_{F,j}$ and $H_{F,j}$ are its horizontal radius and height. With
$\Omega_{\mathrm{hard}}=\bigcup_{j=1}^{N_F}\Omega_{F,j}$, the continuous
collision-avoidance requirement is
\begin{equation}
 \mathbf{r}(t)\notin\Omega_{\mathrm{hard}},\qquad t\in[0,T].
 \label{pr:eq:collision}
\end{equation}

For radar $i$ at $(x_{R,i},y_{R,i},0)$, define the slant range
\begin{equation}
 d_i(\mathbf{p})=\sqrt{(x-x_{R,i})^2+(y-y_{R,i})^2+z^2}.
 \label{pr:eq:slant}
\end{equation}
Its coverage is hemispherical within the nonnegative-altitude workspace.
Let $R_i$ be its outer radius and $R_{\mathrm{in},i}=\eta R_i$, with
$0<\eta<1$. For the unit radar scaling factors used in the experiments,
the local exposure field is
\begin{equation}
 P_{d,i}(\mathbf{p})=
 \begin{cases}
  1, & d_i(\mathbf{p})\le R_{\mathrm{in},i},\\
  \left(\dfrac{R_i-d_i(\mathbf{p})}{R_i-R_{\mathrm{in},i}}\right)^{\nu},
       & R_{\mathrm{in},i}<d_i(\mathbf{p})<R_i,\\
  0, & d_i(\mathbf{p})\ge R_i,
 \end{cases}
 \label{pr:eq:exposurefield}
\end{equation}
where $\nu>0$ controls the decay rate. The experiments use $\eta=0.3$
and $\nu=4$. Overlapping fields are combined by
\begin{equation}
 P_d(\mathbf{p})=\max_{1\le i\le N_R}P_{d,i}(\mathbf{p}),
 \qquad 0\le P_d(\mathbf{p})\le1.
 \label{pr:eq:maxfield}
\end{equation}
This is a prescribed dominant-exposure surrogate, rather than a statistical
fusion formula for the probability of detection by at least one radar.
It does not model radar scan timing, track accumulation or measurement
dependence. Hard zones determine admissibility, while the soft field
contributes to the accumulated exposure objective in
Section~\ref{pr:sec:metrics}.

\subsection{Fixed-wing flight constraints}
\label{pr:sec:constraints}
The model assumes still air and constant aircraft mass. Hence the scalar
path speed $v(t)=\|\dot{\mathbf{r}}(t)\|$ is identified with airspeed.
Feasibility refers to the specified geometric and reduced flight-envelope
constraints; it does not imply validation of a six-degree-of-freedom
aircraft model or a closed-loop tracking controller.

\subsubsection{Mission altitude envelope}
Let $d_s(\mathbf{p})$ and $d_g(\mathbf{p})$ denote the horizontal distances
to the start and goal. For transition distance $D_{\mathrm{tr}}$, define
\begin{equation}
 \alpha_e(\mathbf{p})=\operatorname{clip}
 \left(\frac{d_e(\mathbf{p})}{D_{\mathrm{tr}}},0,1\right),
 \qquad e\in\{s,g\},
 \label{pr:eq:transition}
\end{equation}
and $h_e(\mathbf{p})=z_e+(h_{\mathrm{en}}-z_e)\alpha_e(\mathbf{p})$.
Here $h_{\mathrm{en}}$ is the minimum en-route altitude. The local envelope is
\begin{equation}
 h_{\min}(\mathbf{p})=\min\{h_s(\mathbf{p}),h_g(\mathbf{p}),h_{\mathrm{en}}\},
 \qquad h_{\min}(\mathbf{r}(t))\le z(t)\le h_{\max}.
 \label{pr:eq:altitude}
\end{equation}
This construction retains the endpoint altitudes while imposing the
en-route lower bound outside both transition regions. A sufficient
compatibility condition for the ramp slope is
\begin{equation}
 D_{\mathrm{tr}}\ge
 \max_{e\in\{s,g\}}\frac{\max(h_{\mathrm{en}}-z_e,0)}{\tan\gamma_{\max}}.
 \label{pr:eq:ramp}
\end{equation}
This condition concerns the altitude envelope itself; it does not guarantee
that a feasible route exists in the presence of obstacles. In the present
mission, $z_s=z_g=0$ is an altitude reference for airborne trajectory
segments. Ground roll, take-off dynamics and landing dynamics are not modeled.

\subsubsection{Airspeed constraint}
The aircraft must remain above its stall-speed lower bound. The speed
assignment uses a prescribed cruise ceiling $v_{\mathrm{cr}}$, additionally
bounded by the aircraft maximum speed $v_{\max}$:
\begin{equation}
 v_{\mathrm{stall}}\le v(t)\le v_{\mathrm{cap}},\qquad
 v_{\mathrm{cap}}=\min\{v_{\mathrm{cr}},v_{\max}\}.
 \label{pr:eq:speedcap}
\end{equation}
For the reported configuration, $v_{\mathrm{stall}}=33~\mathrm{m\,s^{-1}}$,
$v_{\mathrm{cr}}=55~\mathrm{m\,s^{-1}}$ and
$v_{\max}=83~\mathrm{m\,s^{-1}}$. The airborne initial condition is
$v(0)=v_{\mathrm{stall}}$. No terminal-speed equality is prescribed;
$v(T)$ follows from the local ceilings and the forward--backward speed
assignment. The current speed assignment seeks the available speed under
these ceilings and does not solve a separate fuel-optimal speed-control problem.

\subsubsection{Turning constraint}
For a coordinated level turn, the bank-limited radius satisfies
\begin{equation}
 R_{\min}(v)=\frac{v^2}{g\tan\phi_{\max}}.
 \label{pr:eq:turnradius}
\end{equation}
The implementation applies this relationship as a curvature-based envelope
test on the three-dimensional path. For a regular spline
$\mathbf{q}(u)$, its spatial curvature is
\begin{equation}
 \kappa(u)=\frac{\|\mathbf{q}'(u)\times\mathbf{q}''(u)\|}
                    {\|\mathbf{q}'(u)\|^3},\qquad
 v_{\mathrm{turn}}(u)=\sqrt{\frac{g\tan\phi_{\max}}{\kappa(u)}}.
 \label{pr:eq:curvature}
\end{equation}
The turning ceiling is interpreted as infinite where $\kappa=0$.
The assigned speed must not exceed both $v_{\mathrm{cap}}$ and
$v_{\mathrm{turn}}$. If $v_{\mathrm{turn}}<v_{\mathrm{stall}}$, the geometry
is rejected; the speed is not reduced below stall to make it acceptable.
Using spatial curvature in this level-turn relationship is a reduced-model
approximation for gently climbing trajectories, not an exact bank-angle
reconstruction for general three-dimensional flight. A constant stall
threshold is used; bank-dependent stall speed and full load-factor dynamics
are outside the present model.

\subsubsection{Climb and vertical-rate constraints}
With horizontal arc length $s_h$, the flight-path angle and vertical rate obey
\begin{equation}
 \gamma=\arctan\left(\frac{\mathrm{d}z}{\mathrm{d}s_h}\right),\qquad
 |\gamma|\le\gamma_{\max},\qquad
 -v_{\mathrm{des},\max}\le\dot z\le v_{\mathrm{asc},\max}.
 \label{pr:eq:vertical}
\end{equation}
The angle bound restricts geometry, while $\dot z=v\sin\gamma$ couples
that geometry to speed. Both are required. The current limits are
$\gamma_{\max}=15^\circ$, $v_{\mathrm{asc},\max}=12~\mathrm{m\,s^{-1}}$
and $v_{\mathrm{des},\max}=15~\mathrm{m\,s^{-1}}$.

\subsubsection{Longitudinal acceleration and numerical feasibility}
The tangential acceleration is bounded by
\begin{equation}
 -a_{\mathrm{dec},\max}\le\dot v(t)\le a_{\mathrm{acc},\max},
 \label{pr:eq:acceleration}
\end{equation}
with $a_{\mathrm{acc},\max}=3~\mathrm{m\,s^{-2}}$ and
$a_{\mathrm{dec},\max}=4~\mathrm{m\,s^{-2}}$ in the experiments.
These bounds are imposed during speed assignment rather than added to the cost.

Search edges are screened for environmental and geometric feasibility.
After spline construction, endpoint retention, terminal horizontal heading,
local deviation, altitude, obstacle clearance and curvature are checked
again, followed by speed-envelope validation. The terminal reference heading
is inferred from the incoming discrete path, rather than supplied as an
independent mission heading. The fine-search corridor restricts graph search;
final spline locality is checked against the input path, and membership in
the coarse corridor is diagnostic by default.

The continuous constraints above specify the intended model. Their
implementation uses finite-resolution sampling and numerical tolerances:
the spline geometry is probed at 4096 parameter values, endpoints are accepted
within $2~\mathrm{m}$, and the terminal-heading tolerance is $12^\circ$.
Therefore, acceptance denotes satisfaction of these numerical tests, not
a formal continuous-domain safety certificate. Numerical settings and
resolution checks must accompany the reported experiments.

\subsection{Time, threat-exposure, and fuel models}
\label{pr:sec:metrics}
The mission metrics are flight time $T$, accumulated exposure $C$, and
fuel consumption $F$. They are evaluated on a common trajectory and with
the same time increments. Let $s\in[0,L]$ denote three-dimensional arc
length along the final curve and let $v(s)>0$ be its assigned path speed.
The associated time law is $\mathrm{d}t=\mathrm{d}s/v(s)$.

\subsubsection{Flight time}
The continuous flight-time functional is
\begin{equation}
 T=\int_0^L\frac{1}{v(s)}\,\mathrm{d}s.
 \label{pr:eq:time}
\end{equation}
For final trajectory samples $\mathbf{q}_i$ with speeds $v_i$, the numerical
evaluator uses the chord length
$\ell_i=\|\mathbf{q}_{i+1}-\mathbf{q}_i\|$ and constant tangential
acceleration on each sampled segment. Consequently,
\begin{equation}
 \Delta t_i=\frac{2\ell_i}{v_i+v_{i+1}},\qquad
 a_i=\frac{v_{i+1}^2-v_i^2}{2\ell_i},\qquad
 T_{\mathrm{num}}=\sum_i\Delta t_i,
 \label{pr:eq:segmenttime}
\end{equation}
for nonzero-length segments. This is exact for the adopted segment motion
model, while the sampled representation approximates the continuous spline.
It avoids replacing the final trajectory speed by one mission-wide constant.

During graph search, the speed profile is not yet available. A straight
edge $e$ of length $\ell_e$ and altitude change $\Delta z_e$ instead uses
\begin{equation}
 \Delta t_e^{\mathrm{search}}=
 \max\left\{\frac{\ell_e}{v_{\mathrm{cr}}},
             \frac{|\Delta z_e|}{v_{z,\max}(e)}\right\},
 \label{pr:eq:searchtime}
\end{equation}
where $v_{z,\max}(e)$ is the ascent or descent limit according to the sign
of $\Delta z_e$; the second term is zero for level flight. Here
$v_{\mathrm{cr}}$ lies within the aircraft speed interval. This edge model
is a search surrogate and does not account for acceleration across graph
vertices. The final metrics are recomputed after geometric refinement and
speed assignment; they need not equal the search-stage estimates.

\subsubsection{Accumulated threat exposure}
The exposure objective is
\begin{equation}
 C=\int_0^T P_d(\mathbf{r}(t))\,\mathrm{d}t
   =\int_0^L\frac{P_d(\mathbf{q}(s))}{v(s)}\,\mathrm{d}s.
 \label{pr:eq:exposure}
\end{equation}
Because $P_d$ is dimensionless and bounded by unity, $C$ has units of
seconds and satisfies $0\le C\le T$. It measures time-weighted exposure,
not mission detection probability, interception probability or expected
loss. The same geometric passage can incur different exposure when the
time spent within the field changes. Thus a length integral of $P_d$ alone
would not represent the objective evaluated here.

\subsubsection{Fuel consumption}
Let $v_h=\sqrt{\dot x^2+\dot y^2}$ and $v_z=\dot z$ be the horizontal and
vertical velocity components. The implemented fuel model uses a
quasi-steady parabolic-drag power surrogate. Air density is approximated by
\begin{equation}
 \rho(z)=\rho_0\exp(-z/H_\rho),\qquad
 \rho_0=1.225~\mathrm{kg\,m^{-3}},\quad H_\rho=8500~\mathrm{m}.
 \label{pr:eq:density}
\end{equation}
This is an exponential atmosphere approximation. It is not the full
piecewise International Standard Atmosphere model. The level-flight
mechanical power surrogate is
\begin{equation}
 P_{\mathrm{lev}}(v_h,z)=
 \frac{1}{2}\rho(z)C_{D0}S v_h^3+
 \frac{2(mg)^2}{\rho(z)v_h S\pi\,\mathrm{AR}\,e_O},
 \label{pr:eq:levelpower}
\end{equation}
where $S$ is wing area, $C_{D0}$ is the zero-lift drag coefficient,
$\mathrm{AR}$ is aspect ratio, and $e_O$ is the Oswald efficiency factor.
The first and second terms represent parasitic and induced power,
respectively. The model uses horizontal speed in these terms; it does not
substitute path speed $v$ for $v_h$ in climbing segments.

For the current configuration, descent energy recovery is zero. With
climb efficiency $\eta_c$ and fuel-to-propulsive efficiency $\eta_p$,
the fuel-equivalent power is
\begin{equation}
 P_f(v_h,v_z,z)=\frac{P_{\mathrm{lev}}(v_h,z)+P_c(v_z)}{\eta_p},
 \qquad
 P_c(v_z)=
 \begin{cases}
  mgv_z/\eta_c, & v_z>\delta_z,\\
  0, & v_z\le\delta_z,
 \end{cases}
 \label{pr:eq:fuelpower}
\end{equation}
where $\delta_z=0.01~\mathrm{m\,s^{-1}}$ is the small vertical-rate
deadband used in evaluation. The routine numerically floors $v_h$ at
$v_{\mathrm{stall}}/2$; this floor is inactive for trajectories satisfying
the stated stall and flight-path-angle bounds. With $H_f$ denoting the fuel
lower heating value,
\begin{equation}
 F=\frac{1}{H_f}\int_0^T
 P_f(v_h(t),v_z(t),z(t))\,\mathrm{d}t,
 \qquad H_f=43\times10^6~\mathrm{J\,kg^{-1}}.
 \label{pr:eq:fuel}
\end{equation}
Thus $F$ is reported in kilograms. Since $P_f$ already includes division
by $\eta_p$, the efficiency must not be applied a second time when
converting energy to fuel mass.

For numerical integration, each sampled segment is subdivided at spatial
fractions $0=\lambda_0<\cdots<\lambda_K=1$. Constant tangential acceleration
gives
\begin{equation}
 v_{ij}=\sqrt{(1-\lambda_j)v_i^2+\lambda_j v_{i+1}^2},\qquad
 \Delta t_{ij}=\frac{2\ell_i(\lambda_{j+1}-\lambda_j)}{v_{ij}+v_{i,j+1}}.
 \label{pr:eq:subtime}
\end{equation}
At subsegment midpoint $\mathbf{q}_{ij}^{m}$, the horizontal and vertical
velocities are calculated from the corresponding displacement divided by
$\Delta t_{ij}$. The three numerical metrics are
\begin{equation}
 \begin{aligned}
 T_{\mathrm{num}}&=\sum_{i,j}\Delta t_{ij},\\
 C_{\mathrm{num}}&=\sum_{i,j}P_d(\mathbf{q}_{ij}^{m})\Delta t_{ij},\\
 F_{\mathrm{num}}&=\frac{1}{H_f}\sum_{i,j}
 P_f(v_{h,ij},v_{z,ij},z_{ij}^{m})\Delta t_{ij}.
 \end{aligned}
 \label{pr:eq:quadrature}
\end{equation}
Spatial midpoint evaluation of the nonlinear field and power terms is a
quadrature approximation. Search edges use the same field and power
functions, with their edge time distributed uniformly over the subsegments.

This fuel model is intended for internally consistent comparisons. It
neglects fuel-induced mass change, bank-dependent induced-drag corrections,
transient engine behavior and an explicit tangential-acceleration power
term. These assumptions limit claims about absolute fuel prediction.

\subsection{Anchor-based normalized mission-oriented objective}
\label{pr:sec:objective}
Time, exposure and fuel have different units and numerical scales. A fixed
reference is therefore constructed from three single-objective searches
on the same coarse graph, with identical mission endpoints, environment,
constraints and search-stage evaluation rules. Denote the resulting paths
by $\Gamma_T$, $\Gamma_C$ and $\Gamma_F$, corresponding to the weight
vectors $(1,0,0)$, $(0,1,0)$ and $(0,0,1)$ in $(T,C,F)$ order. These
are discrete search anchors; they are not independently optimized final
continuous trajectories.

Evaluating all three metrics on each anchor gives the payoff table
\begin{equation}
 A=\begin{bmatrix}
 T^{\mathrm{search}}(\Gamma_T)&C^{\mathrm{search}}(\Gamma_T)&F^{\mathrm{search}}(\Gamma_T)\\
 T^{\mathrm{search}}(\Gamma_C)&C^{\mathrm{search}}(\Gamma_C)&F^{\mathrm{search}}(\Gamma_C)\\
 T^{\mathrm{search}}(\Gamma_F)&C^{\mathrm{search}}(\Gamma_F)&F^{\mathrm{search}}(\Gamma_F)
 \end{bmatrix}.
 \label{pr:eq:payoff}
\end{equation}
For objective $k\in\{T,C,F\}$, let the lower reference be the corresponding
anchor's diagonal entry, and let the upper reference be the column maximum:
\begin{equation}
 b_k=A_{kk},\qquad n_k=\max_{a\in\{T,C,F\}}A_{ak},\qquad
 \Delta_k=n_k-b_k>0.
 \label{pr:eq:anchorscale}
\end{equation}
The upper reference is a payoff-table nadir approximation, not the
maximum over all feasible paths or an exact continuous Pareto nadir.
Degenerate reference ranges must be rejected before normalized planning;
they are not silently replaced by an arbitrary denominator.

For a metric vector $\mathbf{f}=(T,C,F)$, define
\begin{equation}
 \widehat T=\frac{T-b_T}{\Delta_T},\qquad
 \widehat C=\frac{C-b_C}{\Delta_C},\qquad
 \widehat F=\frac{F-b_F}{\Delta_F}.
 \label{pr:eq:normalize}
\end{equation}
The same reference values are retained across weights, coarse and fine
searches, and final-trajectory reporting within a fixed experiment.
The normalized values are not clipped to $[0,1]$: a refined or retimed
trajectory can fall outside the three-anchor reference ranges. This does
not invalidate the normalization, but it precludes interpreting the
reference values as exact bounds on final continuous performance.

Mission preference is expressed through
\begin{equation}
 J=w_T\widehat T+w_C\widehat C+w_F\widehat F,\qquad
 w_T,w_C,w_F\ge0,\quad w_T+w_C+w_F>0.
 \label{pr:eq:objective}
\end{equation}
The coefficients specify relative preferences and need not sum to one.
Multiplication of all weights by the same positive constant preserves
path ordering for a fixed model. During graph search, the additive
contribution of edge $e$ is
\begin{equation}
 c_e=\frac{w_T\Delta T_e}{\Delta_T}
     +\frac{w_C\Delta C_e}{\Delta_C}
     +\frac{w_F\Delta F_e}{\Delta_F}.
 \label{pr:eq:edgeobjective}
\end{equation}
The lower-reference offsets are applied once to the complete path:
\begin{equation}
 J^{\mathrm{search}}(\Gamma_d)=\sum_{e\in\Gamma_d}c_e
 -\left(\frac{w_Tb_T}{\Delta_T}+\frac{w_Cb_C}{\Delta_C}
                         +\frac{w_Fb_F}{\Delta_F}\right).
 \label{pr:eq:offset}
\end{equation}
The bracketed term is path independent. Subtracting the lower references
on every edge would instead introduce a spurious dependence on the number
of edges and would not implement Eq.~\eqref{pr:eq:objective}.

This scalarization directs discrete route selection. Spline candidates
are subsequently ranked with a separate geometric refinement score, and
the final $(T,C,F,J)$ values are recomputed using the assigned speed
profile. The sequential procedure is therefore not a joint global
optimization of route geometry, speed and fuel, nor does it establish
global Pareto optimality of the final trajectories.

The sensitivity experiment fixes $w_C=1$ and evaluates
$w_T\in\{0,0.1,\ldots,1.1\}$ and
$w_F\in\{0,0.1,\ldots,0.5\}$, giving 72 configurations.
The coefficients are swept rather than tuned by an additional optimizer.
For tractability, this sweep uses coarse search followed by the same
types of spline refinement, feasibility checks and final evaluation; it
does not repeat the fine-corridor search for all 72 configurations.
Consequently, it characterizes preference sensitivity of that coarse-search
pipeline. Claims about the full two-stage pipeline require separate
controlled comparisons and representative-case experiments that include
the fine-corridor search.


%% Insert after the COMPLETE Section 2, not inside the fuel model.
\section{Proposed risk-aware two-stage trajectory planning framework}
\label{sec:method}

The framework separates route selection, geometric refinement and speed
assignment while retaining a common mission evaluator. The coarse search
first identifies a route under the normalized time--exposure--fuel objective.
A fine search then resolves its local geometry within a three-dimensional
corridor. Cubic B-splines replace the resulting polyline with differentiable
candidates, which are screened against geometric constraints and ranked
using a dimensionless refinement score. Finally, a forward--backward pass
assigns speeds along the selected curve, and the mission metrics are
recomputed under the resulting motion. This decomposition limits the size
of the search state while making the distinction between search estimates
and final trajectory performance explicit.

%% FIGURE PLACEHOLDER: overall framework; use vector PDF for the final artwork.
%% Show mission/environment -> three raw anchors and shared scales -> coarse A*
%% -> 3-D corridor -> fine A* -> spline candidates/repairs -> geometric checks
%% -> minimum geometric score -> strict speed assignment -> final T, C, F, J.
%% Show rejection exits for search, geometry and speed failures. Add a dashed
%% coarse-to-spline branch labelled "72-case sensitivity protocol".
%% \begin{figure}[t]
%% \centering
%% \includegraphics[width=\linewidth]{figures/framework.pdf}
%% \caption{Planning and evaluation workflow. Mission preferences govern graph
%% search, whereas geometric refinement and speed assignment are subsequent
%% stages with explicit feasibility checks.}
%% \label{fig:framework}
%% \end{figure}

\subsection{Multi-objective risk-aware A* search}
\label{sec:astar}

The search graph represents position on an anisotropic Cartesian lattice
with horizontal and vertical spacings $\delta_{xy}$ and $\delta_z$.
Candidate displacements are drawn from
\begin{equation}
 \mathcal{D}=\{(a\delta_{xy},b\delta_{xy},c\delta_z):
 a,b,c\in\{-1,0,1\},\ (a,b,c)\ne(0,0,0)\}.
 \label{me:eq:neighbors}
\end{equation}
Purely vertical moves are excluded, and a displacement is retained only
if its climb angle satisfies the prescribed bound. The map boundary,
altitude envelope and hard exclusion regions further restrict the
admissible graph. Each retained edge is checked at ten spatial locations,
including its endpoints; the nine intervening subsegments supply the
midpoint quadrature for exposure and fuel. A graph state contains position
only. Heading, curvature and acceleration feasibility are consequently
resolved after search rather than represented as additional state variables.

For a partial path ending at $\mathbf p$, the accumulated search cost is
$G(\mathbf p)=\sum_e c_e$, with $c_e$ defined in
Eq.~\eqref{pr:eq:edgeobjective}. The A* priority is
\begin{equation}
 f(\mathbf p)=G(\mathbf p)+h(\mathbf p),\qquad
 h(\mathbf p)=\frac{w_T}{\Delta_T}
              \frac{\|\mathbf p-\mathbf p_g\|}{v_{\mathrm{cr}}}.
 \label{me:eq:heuristic}
\end{equation}
No heuristic inflation is used. The heuristic retains only a lower bound
on the time contribution: risk and fuel are nonnegative, and the search
time of an edge is at least its length divided by $v_{\mathrm{cr}}$.
For adjacent states $\mathbf p$ and $\mathbf q$, the triangle inequality
therefore gives
\begin{equation}
 h(\mathbf p)-h(\mathbf q)
 \le \frac{w_T\|\mathbf p-\mathbf q\|}{\Delta_Tv_{\mathrm{cr}}}
 \le c_{(\mathbf p,\mathbf q)}.
 \label{me:eq:consistency}
\end{equation}
Thus, the heuristic is consistent for the fixed graph and edge model.
When $w_T=0$, it becomes zero and the primary search reduces to uniform-cost
expansion. Equal primary priorities are ordered using accumulated geometric
length plus Euclidean distance to the goal, providing a deterministic
geometric preference in flat-cost regions. This secondary criterion does
not add a distance penalty to the mission objective.

The constant lower-reference offset in Eq.~\eqref{pr:eq:offset} is restored
when reporting the complete path objective. It is not included repeatedly
in edge costs. The resulting minimum-cost interpretation concerns the
discretized search problem, subject to numerical tolerances and successful
search termination. It does not extend to the subsequently modified and
retimed continuous trajectory.
%% Cite a verified primary A* reference here when finalizing the bibliography.

\subsection{Coarse-to-fine three-dimensional planning}
\label{sec:coarsefine}

A uniform fine lattice over the entire workspace requires many expansions
in regions that are irrelevant to a particular mission preference. The
two-stage search instead uses a coarse path $\Gamma_c$ to restrict the
domain of a second, finer search. For any polyline $\Gamma$, define
\begin{equation}
 d(\mathbf p,\Gamma)=
 \min_{[\mathbf a,\mathbf b]\in\Gamma}\;
 \min_{\tau\in[0,1]}
 \|\mathbf p-\mathbf a-\tau(\mathbf b-\mathbf a)\|.
 \label{me:eq:distance}
\end{equation}
This distance uses the closest point on a segment, rather than the closest
waypoint, and treats horizontal and vertical displacement in the same
three-dimensional metric. The search corridor is
\begin{equation}
 \mathcal C(\Gamma_c,r_c)=
 \{\mathbf p\in\mathcal W:d(\mathbf p,\Gamma_c)\le r_c\}.
 \label{me:eq:corridor}
\end{equation}
Where grid snapping changes an endpoint, the physical endpoint is connected
to the coarse polyline when constructing the corridor reference. The
endpoints used in the present experiments lie on both lattices.

The fine graph retains nodes inside this corridor and samples candidate
edges to prevent a connection between admissible endpoints from leaving
the corridor unchecked. For the reported settings, corridor probes are
spaced no more than approximately $100~\mathrm{m}$ along an edge, with at
least three probes. Altitude and obstacle screening remain active. Both
search stages use the same objective weights and the same anchor-derived
normalization scales; the fine stage does not recompute its own anchors.

The fine result $\Gamma_f$ is obtained by a new graph search, not by
interpolating $\Gamma_c$. Its finer directional and altitude resolution
permits local route changes, but the corridor limits which alternatives
can be explored. The construction therefore trades search coverage for
computational tractability; it does not guarantee equivalence to an
unrestricted fine-grid search. A failed fine search is reported as a
failure rather than relabeling the coarse path as a fine result.

%% FIGURE PLACEHOLDER: coarse/corridor/fine geometry, preferably two panels.
%% Use the SAME actual case: (a) plan view with coarse and fine paths and
%% corridor; (b) 3-D detail or altitude profile showing vertical refinement.
%% Do not draw the 3-D corridor merely as a fixed-width horizontal strip.
%% \begin{figure}[t]
%% \centering
%% \includegraphics[width=\linewidth]{figures/coarse_corridor_fine.pdf}
%% \caption{Coarse-to-fine route refinement within a three-dimensional
%% distance corridor. The corridor restricts the fine search domain.}
%% \label{fig:coarsefine}
%% \end{figure}

\subsection{Cubic B-spline trajectory refinement}
\label{sec:bspline}

The fine polyline satisfies the edge-level search constraints but retains
direction discontinuities at its vertices. Let
$\{\mathbf P_i\}_{i=0}^{n-1}$ denote its ordered vertices after removal of
consecutive duplicates. A cubic parametric B-spline is fitted as
\begin{equation}
 \mathbf q_\sigma(u)=\sum_{j=0}^{K_\sigma-1}
 B_{j,3}(u)\mathbf c_j,\qquad u\in[0,1],
 \label{me:eq:spline}
\end{equation}
where $B_{j,3}$ are cubic basis functions, $\mathbf c_j$ are control points,
and $\sigma$ controls the fit. Chord-length parameterization assigns
$u_i=\sum_{k=0}^{i-1}\|\mathbf P_{k+1}-\mathbf P_k\|/L_{\mathrm{raw}}$,
where $L_{\mathrm{raw}}$ is the polyline length. The initial FITPACK
smoothing fit uses the weighted residual condition
\begin{equation}
 \sum_{i=0}^{n-1}\omega_i^2
       \|\mathbf P_i-\mathbf q_\sigma(u_i)\|^2\le\sigma,
 \qquad
 \omega_i=
 \begin{cases}
 10^5,&i\in\{0,n-1\},\\
 1,&\text{otherwise}.
 \end{cases}
 \label{me:eq:fit}
\end{equation}
Large endpoint weights promote endpoint retention, which is subsequently
verified with an explicit tolerance. The smoothing condition controls the
initial fit; it is not the final mission objective and is not assumed to
remain an exact residual bound after geometric repair.

To explore the balance between polyline fidelity and curvature reduction,
the smoothing budgets are scaled to the input resolution:
\begin{equation}
 \sigma_j=n(\delta_{xy}\beta_j)^2.
 \label{me:eq:budgets}
\end{equation}
The fixed set of 21 dimensionless factors is
\begin{equation}
 \begin{split}
 \mathcal B=\{&0.02,0.03,0.04,0.05,0.0625,0.075,0.09,\\
              &0.105,0.125,0.15,0.18,0.22,0.25,0.35,\\
              &0.50,0.75,1.0,1.5,2.0,3.0,4.0\}.
 \end{split}
 \label{me:eq:budgetset}
\end{equation}
Every budget is examined before selecting a candidate. Small budgets
retain the lattice geometry more closely; larger budgets allow greater
deformation but must still satisfy the local-deviation and flight-envelope
tests. The input is $\Gamma_f$ for the complete two-stage method and
$\Gamma_c$ for the separate coarse-search sensitivity protocol.

Geometric repair addresses constraint violations introduced by smoothing.
Altitude repair raises deficient sampled altitudes relative to the spatial
altitude floor and refits the corrected points. Climb repair limits the
vertical slope using margins below $\tan\gamma_{\max}$ while retaining
the sampled horizontal coordinates before refitting. These operations
generate revised candidates, each of
which is rechecked against the full geometric test set. Repair is a
feasibility procedure and need not reduce the mission objective.

The terminal heading is inferred from the incoming raw route over a
lookback length of $\max(5\delta_{xy},1000~\mathrm{m})$, using the horizontal
chord from the selected upstream vertex to the goal. If the initial
candidate pass finds no geometrically feasible curve and includes a
heading violation, a second pass aligns the final control-polygon edge
with this reference and reevaluates all smoothing budgets. The adjustment
retains the endpoint coordinates and control-point altitudes. The method
thus preserves the approach direction implied by the route; it does not
enforce an independently specified terminal mission heading.

For regular curves with simple interior knots, cubic B-splines provide
$C^2$ geometric continuity and permit curvature evaluation from analytic
derivatives. This property removes polyline corners, but does not by itself
establish feasibility or imply continuous jerk in the time-parameterized
motion. Those issues require the subsequent geometric and speed checks.
%% Cite the verified FITPACK/Dierckx reference used by the installed SciPy
%% implementation when finalizing the bibliography.

\subsection{Geometric feasibility validation and candidate selection}
\label{sec:trajectory_validation}

Candidate curves are evaluated at $M=4096$ uniformly spaced parameter
values. Their analytic derivatives give spatial curvature through
Eq.~\eqref{pr:eq:curvature} and the local flight-path angle
\begin{equation}
 \gamma(u)=\arctan\!left(
 \frac{|q'_z(u)|}{\sqrt{q'_x(u)^2+q'_y(u)^2}}\right).
 \label{me:eq:splineangle}
\end{equation}
Finite coordinates and nonzero tangents are required. The geometric
turning-radius threshold is
\begin{equation}
 R_{\mathrm{ref}}=
 \max\left\{R_{\mathrm{floor}},
       \frac{v_{\mathrm{stall}}^2}{g\tan\phi_{\max}}\right\}.
 \label{me:eq:radius}
\end{equation}
Requiring $\kappa\le1/R_{\mathrm{ref}}$ prevents acceptance of a turn
that already exceeds the bank envelope at stall speed. It does not require
the curve to support cruise speed everywhere, because speed can subsequently
decrease in curved sections.

A candidate must also satisfy the altitude envelope, avoid the hard
exclusion regions, retain both endpoints within $2~\mathrm{m}$, and match
the terminal horizontal heading within $12^\circ$. Its maximum distance
from the input polyline must satisfy
\begin{equation}
 d(\mathbf q_\sigma(u),\Gamma_{\mathrm{raw}})\le r_s
 \label{me:eq:locality}
\end{equation}
at the validation samples. Here $r_s$ is a refinement locality limit,
distinct from the coarse-search corridor radius $r_c$. In the default
configuration, membership of the final spline in the coarse corridor is
recorded diagnostically, whereas locality relative to the raw input curve
is enforced. Obstacles and the altitude envelope remain hard requirements.

Among the geometrically admissible candidates, selection balances bending,
curvature variation, deviation and length. Let $s$ denote arc length,
$L_\sigma$ the candidate length, and $\kappa_s=\mathrm d\kappa/\mathrm ds$.
The dimensionless components are
\begin{align}
 Q_b&=\frac{1}{L_\sigma}
       \int_0^{L_\sigma}(R_{\mathrm{ref}}\kappa)^2\,\mathrm ds,
       \label{me:eq:bending}\\
 Q_v&=R_{\mathrm{ref}}^3
       \int_0^{L_\sigma}\kappa_s^2\,\mathrm ds,
       \label{me:eq:variation}\\
 Q_d&=\frac{1}{r_s}
       \sqrt{\frac{1}{M}\sum_{m=0}^{M-1}
       d(\mathbf q_\sigma(u_m),\Gamma_{\mathrm{raw}})^2},
       \qquad u_m=\frac{m}{M-1},
       \label{me:eq:deviation}\\
 Q_l&=\frac{L_\sigma}{L_{\mathrm{raw}}}.
       \label{me:eq:length}
\end{align}
The integrals use the cumulative chord-length coordinate of the dense
samples and trapezoidal quadrature; $\kappa_s$ is estimated numerically
on that coordinate. In contrast, $Q_d$ is the root-mean-square deviation
over uniform parameter samples, not an arc-length-weighted integral.
The curvature-variation term penalizes spatial changes in curvature;
it is not a roll-rate constraint.

The refinement score and selection rule are
\begin{equation}
 J_s=0.4Q_b+0.3Q_v+0.2Q_d+0.1Q_l,\qquad
 \mathbf q^*=\mathop{\arg\min}_{\mathbf q\in\mathcal Q_{\mathrm{geom}}}
 J_s(\mathbf q),
 \label{me:eq:selection}
\end{equation}
where $\mathcal Q_{\mathrm{geom}}$ contains the admitted candidates from
the applicable fitting and repair pass. The four coefficients are fixed
geometric design settings, independent of $(w_T,w_C,w_F)$. Selection
enumerates the candidate set rather than stopping at the first feasible
smoothing budget. It is a finite candidate comparison, not a global
optimization over all spline control points.

If no geometrically admissible candidate exists, refinement fails. A
diagnostic fallback curve is not accepted as a mission trajectory. The
selected curve must next pass strict speed validation; the present
procedure does not rerank alternative candidates following a speed-stage
failure. Final acceptance therefore requires both stages to succeed.
Because geometric tests use finite sampling, acceptance has the numerical
meaning stated in Section~\ref{pr:sec:constraints}.

\subsection{Adaptive speed planning and final mission evaluation}
\label{sec:speed}

The selected spline is resampled at approximately uniform arc-length
spacing. A dense set of 8000 parameter probes supplies the cumulative
length map, which is inverted to obtain the final samples
$\{\mathbf q_i\}_{i=0}^{N}$. Their chord lengths are
$\ell_i=\|\mathbf q_{i+1}-\mathbf q_i\|$, and $\kappa_i$ is evaluated
from the spline derivatives at the corresponding parameters. The target
spacing is $100~\mathrm{m}$ in both experimental protocols.

The vertical-rate ceiling of a nondegenerate segment is
\begin{equation}
 \bar v_{z,i}=\begin{cases}
 v_{\mathrm{asc},\max}\ell_i/|z_{i+1}-z_i|,&z_{i+1}>z_i,\\
 v_{\mathrm{des},\max}\ell_i/|z_{i+1}-z_i|,&z_{i+1}<z_i,\\
 +\infty,&z_{i+1}=z_i.
 \end{cases}
 \label{me:eq:verticalcap}
\end{equation}
This ceiling is assigned to both endpoints of the segment. At node $i$,
the admissible speed ceiling is
\begin{equation}
 \bar v_i=\min\left\{
 v_{\mathrm{cr}},v_{\max},
 \sqrt{\frac{g\tan\phi_{\max}}{\kappa_i}},
 \min_{e\in\mathcal E_i}\bar v_{z,e}\right\},
 \label{me:eq:nodecap}
\end{equation}
where $\mathcal E_i$ comprises adjacent segments and the curvature term is
infinite on a straight segment. A ceiling below $v_{\mathrm{stall}}$
causes rejection, rather than a speed assignment below stall.

With $v_0^+=v_{\mathrm{stall}}$, the forward pass propagates the acceleration
bound:
\begin{equation}
 v_{i+1}^+=\min\left\{\bar v_{i+1},
       \sqrt{(v_i^+)^2+2a_{\mathrm{acc},\max}\ell_i}\right\},
 \qquad i=0,\ldots,N-1.
 \label{me:eq:forward}
\end{equation}
The backward pass anticipates downstream speed reductions:
\begin{equation}
 v_N=v_N^+,\qquad
 v_i=\min\left\{v_i^+,
       \sqrt{v_{i+1}^2+2a_{\mathrm{dec},\max}\ell_i}\right\},
 \qquad i=N-1,\ldots,0.
 \label{me:eq:backward}
\end{equation}
No terminal-speed equality is introduced. Applying a segment's vertical
ceiling at both ends bounds its largest vertical speed under the adopted
constant-tangential-acceleration segment model. The final validation checks
speed, turning, climb angle, vertical rate and tangential acceleration;
failure prevents the trajectory from being reported as feasible.

The accepted speed sequence supplies segment times through
Eq.~\eqref{pr:eq:segmenttime}. Time, exposure and fuel are then evaluated
with the common subsegment quadrature in Eq.~\eqref{pr:eq:quadrature},
including renewed altitude and collision screening, and normalized using
the unchanged anchor scales. This last evaluation is essential because
spline deformation changes the spatial threat field sampled by the
vehicle, while retiming changes both residence time and fuel expenditure.
Consequently, a lower geometric refinement score does not by itself imply
a lower final mission objective. The reported output is the accepted
geometry, its assigned speed profile and the recomputed $(T,C,F,J)$;
search-stage estimates are retained separately for diagnosis and ablation.

%% FIGURE PLACEHOLDER: use measured outputs in Section 4.2, not synthetic data.
%% Suggested panels: local raw/spline geometry; analytic spline curvature
%% with 1/R_ref; speed with stall/local ceilings; climb and vertical-rate
%% profiles with limits. Raw-polyline corner curvature is diagnostic only.


\section{Numerical results and discussion}
\label{sec:results}

%% 第四章执行总顺序（以下 %% 均为写作/实验备注，不进入编译正文）：
%% 1. 冻结并备份当前代码、baseline、72 组结果及运行环境，先验收数据，再画图。
%% 2. 完成 4.1 参数与环境图；不运行新搜索即可完成这一部分。
%% 3. 对现有 sweep 做 4.3，得到确定的代表权重；完成选择后再做 4.2/4.4。
%% 4. 单独运行时间单目标完整流程，以及所选代表权重的完整两阶段流程。
%%    若代表权重重复，只运行一次，并如实报告代表规则得到同一个解。
%% 5. 用每次 main 的粗/细原始/最终路径完成 4.2，用代表案例完成 4.4。
%% 6. 最后写第四章开头的实验概览、摘要结果句和结论，不提前填提升百分比。
%% 已完成且签名一致的 72 组不用为了画图再次运行。不要使用 --fresh 作为画图步骤。
%% 下面正文只给出代码可核实的设置；4.2--4.4 的结果叙述等待实测数据。

\subsection{Simulation setup}
\label{sec:setup}

The simulation uses a $100\times100~\mathrm{km}$ horizontal workspace
containing five ground-based radar sites and five associated cylindrical
exclusion zones. The radar layout is generated with seed 42 and held fixed
throughout the preference comparison. Each radar has a hemispherical
coverage radius of $15~\mathrm{km}$, while its colocated hard zone has a
horizontal radius of $6~\mathrm{km}$ and a height of $5~\mathrm{km}$.
The power-law exposure model uses an inner-radius ratio of 0.3 and an
exponent of 4, with overlapping fields combined by their maximum.
These settings define a controlled synthetic environment; they are not
calibrated measurements of a particular radar or aircraft system.

The mission endpoints are $(2,2,0)$ and $(98,98,0)~\mathrm{km}$.
The en-route altitude interval is $1$--$10~\mathrm{km}$, with a
$5~\mathrm{km}$ horizontal transition distance at each endpoint.
The zero endpoint altitude is the reference used for the airborne
trajectory model, as described in Section~\ref{pr:sec:constraints}.
The vehicle and energy parameters are listed in Table~\ref{tab:aircraft};
the numerical planning settings are summarized in
Table~\ref{tab:planning}.

\begin{table}[t]
 \centering
 \caption{Parameters of the reduced fixed-wing model.}
 \label{tab:aircraft}
 \small
 \begin{tabular}{ll}
 \hline
 Parameter & Value \\
 \hline
 Constant mass & $1100~\mathrm{kg}$ \\
 Wing reference area & $11.5~\mathrm{m^2}$ \\
 Gravitational acceleration & $9.81~\mathrm{m\,s^{-2}}$ \\
 Stall / cruise / aircraft maximum speed & $33 / 55 / 83~\mathrm{m\,s^{-1}}$ \\
 Maximum ascent / descent rate & $12 / 15~\mathrm{m\,s^{-1}}$ \\
 Tangential acceleration / deceleration bound & $3 / 4~\mathrm{m\,s^{-2}}$ \\
 Maximum bank / flight-path angle & $30^\circ / 15^\circ$ \\
 Structural turning-radius floor & $50~\mathrm{m}$ \\
 Zero-lift drag coefficient / aspect ratio & $0.025 / 18$ \\
 Oswald efficiency & $0.85$ \\
 Climb / overall propulsive efficiency & $0.80 / 0.30$ \\
 Descent energy recovery fraction & $0$ \\
 Fuel lower heating value & $43~\mathrm{MJ\,kg^{-1}}$ \\
 Sea-level density / density scale height & $1.225~\mathrm{kg\,m^{-3}} / 8500~\mathrm{m}$ \\
 \hline
 \end{tabular}
\end{table}

\begin{table}[t]
 \centering
 \caption{Search and trajectory-evaluation settings. The sweep and the
 complete two-stage method use different input resolutions and locality limits.}
 \label{tab:planning}
 \small
 \begin{tabular}{ll}
 \hline
 Parameter & Value \\
 \hline
 Coarse horizontal / vertical spacing & $2000 / 500~\mathrm{m}$ \\
 Fine horizontal / vertical spacing & $200 / 50~\mathrm{m}$ \\
 Three-dimensional fine-search corridor radius & $600~\mathrm{m}$ \\
 Spline locality limit, fine input & $200~\mathrm{m}$ \\
 Spline locality limit, coarse sweep input & $3015.90~\mathrm{m}$ \\
 Search edge probes / quadrature subsegments & $10 / 9$ \\
 A* heuristic inflation factor & $1.0$ \\
 Initial smoothing budgets & $21$ \\
 Spline geometric validation samples & $4096$ \\
 Arc-length inversion probes & $8000$ \\
 Final target sampling distance & $100~\mathrm{m}$ \\
 Endpoint / terminal-heading tolerance & $2~\mathrm{m} / 12^\circ$ \\
 Initial speed & $33~\mathrm{m\,s^{-1}}$ \\
 Terminal speed & No prescribed equality \\
 \hline
 \end{tabular}
\end{table}

Two protocols address different experimental questions. The complete
two-stage protocol comprises coarse search, corridor-restricted fine search,
spline refinement and strict speed validation. It supports the controlled
comparisons and representative-case assessment. The sensitivity protocol
uses coarse search followed by refinement and speed validation for each
of the 72 weight configurations specified in
Section~\ref{pr:sec:objective}. Its locality limit is scaled to the coarse
cell diagonal,
\begin{equation}
 r_s^{\mathrm{sweep}}=1.05\sqrt{2(2000)^2+(500)^2}
                    \simeq3015.90~\mathrm{m},
 \label{ex:eq:sweeptrust}
\end{equation}
whereas the complete method uses $r_s=200~\mathrm{m}$ relative to the fine
polyline. The two protocols share the threat and aircraft models, objective
scales, candidate-ranking rule and final metric definitions. Their
different search domains and locality limits are retained explicitly
when interpreting the results.

Search-stage and final-trajectory quantities are distinguished throughout
the analysis. Flight time $T$, exposure $C$ and fuel $F$ are reported in
seconds, exposure-seconds and kilograms, respectively; path length is
reported in kilometres. The exposure integral measures time-weighted
presence in the prescribed field, not a mission-level detection probability.
Cross-preference comparisons use the physical vector $(T,C,F)$: a value
of $J$ computed with one weight vector cannot directly rank a trajectory
scored with another. For a controlled comparison at fixed weights, the
shared anchor normalization provides the common scalar objective.

%% 4.1 操作与核对：
%% A. 环境来自 radar_environment_5zones.py::generate(seed=42)，正文使用 main.py
%%    实际端点 (2,2,0)/(98,98,0) km，不用环境类内部默认的 (5,5)/(95,95)。
%% B. 从运行缓存/原始环境导出五站坐标；归档精确值，不用截图反读坐标。
%% C. 填真实实验机器 CPU、RAM、操作系统、Python/NumPy/SciPy 版本。
%%    本次写作机器的软件版本不能自动充当先前实验的版本。
%% D. 归档 baseline.pkl 中的三个 anchor 原始 T/C/F 与六个尺度参考值；
%%    三条 anchor 未做最终样条处理，不混进“最终轨迹”比较表。
%% E. 计时应分别报告 anchor（一次性）、粗搜索、细搜索、样条、速度/评价。
%%    不把绘图/Excel 导出时间混进搜索时间，不把缓存加载时间当重新计算时间。
%%    只有实际补做重复测量后才能写重复次数、均值/标准差；同图重复计时
%%    反映计算波动，不是多环境统计证据。未测步骤写未记录，不填 0。
%% F. preflight_check.py 可作运行前检查；程序测试不能代替实验数据验收。
%% G. 建议最小数值核查：对选定轨迹另做更密集几何/积分采样，比较最小净空、
%%    最大曲率、最大角度和 T/C/F 的变化；固定几何时无需重跑 A* 或 72 组。
%%    100 -> 50 m 若重新分配速度，属于“离散化+速度策略”联合变化，须单列；
%%    不把它说成仅积分误差。未完成上述核查前不宣称数值收敛或连续安全认证。

%% FIGURE PLACEHOLDER: environment overview; let LaTeX assign figure numbers.
%% Panels: (a) 3-D hemispheres and hard cylinders with mission endpoints;
%% (b) exposure slices at stated altitudes with shared [0,1] color scale.
%% Hard-zone disk only appears when the slice altitude is <= 5 km.
%% Use physical axis units and distinguish hard zones from traversable exposure.
%% \begin{figure}[t]
%% \centering
%% \includegraphics[width=\linewidth]{figures/threat_environment.pdf}
%% \caption{Synthetic mission environment with five radar fields and colocated
%% hard exclusion zones. Soft exposure remains traversable and contributes
%% to the time-integrated objective.}
%% \label{fig:environment}
%% \end{figure}

\subsection{Overall performance and controlled ablation}
\label{sec:ablation}

%% 本节回答三个不同问题：任务代价是否改变选路；细网格改变了哪些搜索几何；
%% 样条与速度规划如何得到被数值验收的飞行轨迹。不要把三种变化混成单一增益。
%% 建议总表：行=Time-oriented final / Risk-aware final / Coarse raw /
%% Fine raw（后两行按同一风险权重）；列=表示阶段、权重、L、T、C、F、J、
%% 节点数、分阶段耗时。Raw 行的指标加 search 上标，不与 final 行算“总提升”。
%% 几何/动力学指标另表或配图，避免一张表过宽。表头明确单位及检查口径。

\subsubsection{Time-oriented A* versus risk-aware planning}
\label{sec:objective_comparison}

%% 研究问题：在相同离散结构与后处理下，加入暴露/燃油偏好是否改变路线，
%% 带来多大的暴露--时间--燃油权衡，而不是证明 A* 被一种新搜索机制击败。
%% 步骤 1：定义 baseline 权重 (w_T,w_C,w_F)=(1,0,0)。这是时间单目标 A*；
%%         因垂直率限时，不能无条件写成 shortest-distance A*。
%% 步骤 2：固定一组公开选择规则得到的 risk-aware 权重（例如 4.3 的平衡解）；
%%         不挑一个只让结果最好看的权重；main 默认 (0.3,1,0.2) 不是自动平衡解。
%% 步骤 3：两组均用完整两阶段流程、同一环境/端点/网格/半径/飞机/归一化/
%%         平滑与速度算法。路线不同导致 corridor 不同，这是目标改变的后果。
%%         该实验评估完整流程的偏好效果，不是严格锁定同一个 fine graph。
%% 步骤 4：每次运行先归档 planned_path.pkl、planned_path_coarse.pkl、baseline、
%%         完整日志及参数/代码快照；main 默认文件名会被下一次运行覆盖。
%% 步骤 5：仅在双方最终可行时，列最终 L、T、C、F、最高高度及计算耗时。
%%         想比较 J 时，用同一个声明的参考权重重新评价两条最终轨迹；
%%         不比较各自不同权重下的 J，且 J 可为负，不宜笼统算百分比下降。
%% 步骤 6：物理指标百分比采用 100*(new-base)/base，明确负值表示下降；
%%         base=0（尤其 C=0）时报告绝对差，不加任意 epsilon 制造百分数。
%% 图：同一坐标/同一高度风险说明下叠加两条三维路径和俯视路径，补高度剖面。
%% 段落顺序：共同设置 -> 数字差异 -> 通过/绕过/爬升避让的几何证据 -> 代价。
%% 结论边界：可说“在该场景和偏好下”，不能推论所有环境都降低 T、C、F。
%% 若某方案失败，报告失败阶段和条件，不把缺失值补成零或隐去该工况。
%% \begin{figure}[t]
%% \centering
%% \includegraphics[width=\linewidth]{figures/time_vs_risk_aware.pdf}
%% \caption{Final trajectories obtained with time-oriented and risk-aware
%% objectives under the common two-stage planning and validation protocol.}
%% \label{fig:objective_comparison}
%% \end{figure}

\subsubsection{Coarse versus fine planning}
\label{sec:resolution_comparison}

%% 研究问题：在同一次两阶段运行中，细网格对粗路线做了哪些局部调整？
%% 数据：同一次归档的 planned_path_coarse.pkl['path','stats'] 与
%%       planned_path.pkl['path_raw','stats_raw']，不是后者的最终 path/stats。
%% 步骤 1：核对权重、baseline 和环境一致；对齐原始粗/细两条路径。
%% 步骤 2：列搜索口径 L、T_search、C_search、F_search、J_search、扩展节点、
%%         coarse_time、fine_time；细阶段时间不等于完整两阶段总时间。
%% 步骤 3：画三维/局部放大与高度剖面，解释网格转向和高度选择改变的位置。
%% 步骤 4：检查细路径到粗中心线的三维距离及 600 m 限制；不可拿二维投影
%%         距离替代真实三维距离。最终 spline 可能出 coarse corridor，另属后处理。
%% 步骤 5：报告改善、不变或变差；原始粗边可能不属于 corridor 中的细图边，
%%         因此没有先验“细网格目标必定更优”的保证。
%% 图可直接复用 Section 3 的 coarse_corridor_fine 图，在此引用，避免重复图片。
%% 不可把该对比写成“相对全域细网格节省 X%”：那需要另运行全域细网格基线。
%% 不可用 coarse raw vs fine final 把细化、平滑、变速三种效果合并称网格收益。

\subsubsection{Fine raw path versus B-spline trajectory}
\label{sec:refinement_comparison}

%% 研究问题：连续化改动多少几何，经过哪些验证，可执行性代价是什么？
%% 数据：同一个 planned_path.pkl 中 path_raw、path、smooth_result、
%%       curvature_smooth、speed_profile、accel_profile、stats_raw、stats。
%% 步骤 1：画原始/最终局部转弯及高度剖面；给 L、最大/RMS 局部偏离、端点误差、
%%         末端航向误差，说明是否满足 200 m、2 m、12 deg 等对应阈值。
%% 步骤 2：最终曲率必须用保存的解析 curvature_smooth 或 tck 导数；原始折线
%%         顶点曲率不连续，三点估计仅为网格诊断，不能宣布其连续最大曲率有限。
%% 步骤 3：画最终曲率与 1/R_ref、速度与失速/局部上限、切向加减速度与 -4/+3、
%%         爬升角与 15 deg、垂直率与 -15/+12。用带符号加速度区分加速与减速。
%%         垂直率验收包含段内最大值，不仅看平均 dz/dt 曲线。
%% 步骤 4：报告最小硬区净空/碰撞检测、最低高度包络余量、全部几何验收状态；
%%         一张看起来平滑的图不能代替这些检查。有限采样不等于连续证明。
%% 步骤 5：列候选数、可行候选数、所选 sigma、J_s 分量、是否航向修复；
%%         21/42 是平滑预算尝试数，内部修复可能多次验算，不能当总验证调用数。
%% 步骤 6：T_search/C_search/F_search 与 final T/C/F 分列，说明差别含几何与
%%         速度变化；不能把原始折线称为已满足固定翼转弯约束的可飞行对照。
%%         若要单独归因速度规划收益，需追加同一几何的明确可行速度对照。
%% 段落顺序：折线局限 -> 局部几何修整 -> 约束和速度证据 -> 最终任务代价。
%% \begin{figure}[t]
%% \centering
%% \includegraphics[width=\linewidth]{figures/refinement_feasibility.pdf}
%% \caption{Geometric refinement and speed-envelope validation for the same
%% fine-search route. Constraint limits accompany the final trajectory profiles.}
%% \label{fig:refinement}
%% \end{figure}

\subsection{Mission-preference sensitivity}
\label{sec:sensitivity}

%% 本节仅使用 coarse-search + refinement + speed 的 72 组协议，不能写成
%% 72 组都执行了 corridor fine search，也不把重复权重当独立随机实验。
%% 数据：output_5zones/sweep_2d_time_fuel_v5_3_1.xlsx 的“全部数据”工作表，
%%       与 sweep_cases_v5_3_1 中匹配的逐工况缓存；入口 analysis_data_loader.py。
%% 第一步【验收】：确认 12*6 唯一组合、w_C=1、schema/平滑版本/run_signature
%% 一致且匹配当前基准和代码、所有成功标记和物理指标有效。不能把失败行删除
%% 后仍声称完成 72 个有效工况。未完成时先报告进度，不制作正式完整热力图。
%% 第二步【三目标响应】：运行 analyze_weight_sensitivity.py，画 T/C/F 三幅
%% 热力图，x=w_T，y=w_F，使用相同网格顺序并标单位，保留零权重行列。
%% (w_T,w_F)=(0,0) 是纯暴露目标，因为 w_C=1，并非“全零权重无目标”。
%% 先报各指标实测范围和位置，再解释横向/纵向变化；不预写严格单调关系。
%% 相同权重改变可能跨离散路线分支；阶梯、平台与跳变应联系路线/高度证据。
%% 第三步【控制变量切片】：固定 w_F 比较递增 w_T，固定 w_T 比较递增 w_F。
%% 选择提前说明的切片（例如边界及中间值），不要只展示有利曲线。
%% 绝对数值优先；可给相邻格差值，跳变点回查原始路径和样条候选变化。
%% 第四步【权衡散点】：用 analyze_core_tradeoffs_v2.py 的现有输出辅助，
%% 画 T-C 散点以 F 着色，必要时补 C-F/T-F。检查每幅图采用物理量还是批内缩放。
%% 如报告非支配集：对完整实测 T/C/F 做三目标逐对支配判断（三项<=且至少一项<），
%% 去重规则/容差公开；称“本次采样集的非支配解”，不称真实全局 Pareto 前沿。
%% 二维投影中的低边界不能直接等同于三目标非支配；算法也未证明扫全非凸前沿。
%% 第五步【高度机制】：用 analyze_altitude_strategy_v2.py 和个别缓存复查；
%% 最大高度只说明峰值，不等于全程巡航高度。联系上升耗能、密度、几何距离及
%% 斜距暴露解释实际模式，不能先假定“越高必然越省油/越安全”。
%% 第六步【代表选择】：运行 select_representative_weights.py，公开以下规则：
%% 最小暴露排序 C,T,F,w_T,w_F；最短时间排序 T,C,F,w_T,w_F；
%% 最少燃油排序 F,C,T,w_T,w_F；平衡解按批内 T/C/F min-max 归一化后的
%% 欧氏理想点距离排序，再按 T,C,F,w_T,w_F 打破并列。常数列贡献为 0。
%% 这是代表筛选用批内归一化，不是 Section 2 的 anchor 归一化。
%% 指纹/平台邻居数仅为诊断，不参与选择；相同舍入指标不能证明路径完全相同。
%% 如报告权重平台“稳健”，仅指本场景离散权重邻域响应，不指风/雷达误差鲁棒性。
%% 建议正文三段：总体响应 -> 权衡/路线机制 -> 代表选择与局限。
%% \begin{figure}[t]
%% \centering
%% \includegraphics[width=\linewidth]{figures/weight_sensitivity.pdf}
%% \caption{Time, exposure and fuel across the 72 preference configurations
%% evaluated using coarse search followed by trajectory refinement and retiming.}
%% \label{fig:sensitivity}
%% \end{figure}
%% \begin{figure}[t]
%% \centering
%% \includegraphics[width=\linewidth]{figures/mission_tradeoffs.pdf}
%% \caption{Observed mission trade-offs within the evaluated preference set.
%% Marked cases are selected using the stated physical-objective rules.}
%% \label{fig:tradeoffs}
%% \end{figure}

\subsection{Representative mission scenarios}
\label{sec:representative}

%% 目的：检验 coarse sweep 选出的任务偏好在完整两阶段流程中表现如何。
%% 这是同一环境的不同任务偏好案例，不是三个独立环境；标题亦可改为
%% Representative mission preferences，以避免读者误解。
%% 步骤 1：读取 representative_weights.csv，优先选最小暴露/平衡/最短时间。
%%         若燃油权衡是核心贡献，补最少燃油，或在表中列出该解并解释重合。
%%         选择规则固定；不同规则可选同一权重，不强行造三个不同代表。
%% 步骤 2：先保存当前 main.py 配置与现有结果，设置每个已选权重并运行完整
%%         main.py，分别归档源码配置、两种 pkl、日志、环境和 baseline。
%%         更改 main 权重会改变 main 结果签名；跨案例比较应使用各自归档
%%         配置核对，不能手改缓存签名骗过当前加载器。不要改 sweep 源码。
%% 步骤 3：每个案例核对 final success、包络/转弯/速度/加速/航向验证；失败
%%         如实报告，不悄悄换成 sweep 结果冒充 full-pipeline final。
%% 步骤 4：表列选择规则、(w_T,w_C,w_F)、最终 L/T/C/F、最高高度、最小半径、
%%         最大爬升角、航向误差、计算耗时。J 仅在同一参考权重下才可横向比较。
%% 步骤 5：叠加三维/俯视轨迹、统一横轴的高度剖面，解释绕飞或穿越软暴露区。
%%         暴露累计曲线须使用与最终 evaluator 相同的变速积分时间；终值应为 C。
%%         不要用旧恒速积分函数生成“累计暴露曲线”再与变速最终 C 混用。
%% 步骤 6：若对比 coarse-screening vs full-pipeline，逐项报告 T/C/F 变化。
%%         这同时改变 fine search 与 locality 半径（约3016 m -> 200 m），
%%         只能解释为协议差异，不能单独归因于 finer grid。
%% 步骤 7：说明代表性权重是在 coarse 协议中选出，未证明 full-pipeline 下
%%         仍是全72组最小暴露/最快/最省油。可称“由最小暴露规则选出的偏好”。
%% 段落顺序：选择与执行协议 -> 每个案例的空间/速度行为 -> 实测代价 -> 适用边界。
%% 不同权重的效果可能在完整流程中收敛、反转或失败；以结果为准。
%% \begin{figure}[t]
%% \centering
%% \includegraphics[width=\linewidth]{figures/representative_missions.pdf}
%% \caption{Complete two-stage trajectories for preferences selected from the
%% coarse-search sensitivity experiment, shown in the common threat environment.}
%% \label{fig:representatives}
%% \end{figure}

%% 第四章收尾：仅在完成数据后写 1 段综合讨论。分别回答偏好作用、细化作用、
%% 可行性证据；列明计算开销及代价，而非只列优势。承认单一合成布局、简化
%% 气动/能耗与静态威胁、有限候选和采样的适用范围。无需为写作而重跑全部72组。
%% 若要声称跨场景泛化、真实飞行鲁棒或统计显著，必须另设计相应实验，不能
%% 用本节同一环境的权重扫频替代。这些尚未做的工作可放未来工作。


\section{Conclusions}
\label{sec:conclusions}
%% 完成第四章后写：方法解决的问题 -> 经实测支持的关键结果 -> 代价和边界。
%% 不声称全局最优、真实飞行验证、跨场景鲁棒或未进行的显著性检验。
%% 补充未来工作时与实际局限对应，避免用展望代替本稿必须提供的证据。

%% BIBLIOGRAPHY PLACEHOLDER: replace with verified, relevant primary references.
%% The template's example appendix and Lamport example citation were removed.
%% Add the A*, FITPACK/B-spline, flight-envelope and related-work references
%% after checking their metadata and the claims for which they are cited.
%% \bibliographystyle{elsarticle-num}
%% \bibliography{references}

\end{document}
